% \documentclass[12pt, draft]{article} % not showing the graphs
\documentclass[11pt]{article}
\usepackage[utf8]{inputenc}
\usepackage[a4paper, total={6in, 8in}]{geometry}
\usepackage[T1]{fontenc}
\usepackage{lipsum}
\usepackage{setspace}
% \onehalfspacing % Set line spacing to 1.5

% \usepackage{biblatex}
\usepackage[backend=biber,style=authoryear]{biblatex}
\usepackage{caption}
\usepackage{url}
\usepackage{hyperref}
\usepackage{graphics}
\usepackage{bbm}
\usepackage{derivative}
\usepackage{commath, csquotes}
\usepackage[T1]{fontenc}
%\usepackage[showframe]{geometry}
\addbibresource{references.bib}

\usepackage{hyperref}
\usepackage[T1]{fontenc}
\usepackage{tikz}
\usepackage{xurl}
\usepackage{times}
\usepackage{textcomp}
\usepackage{istgame}
\usepackage{latexsym,xcolor,multicol,booktabs,calligra}
\usepackage{amsmath,amssymb,BOONDOX-cal,bm}	
\usepackage{graphicx,pstricks,stackengine}     
\usepackage{comment}     

\usepackage{tabls}
\usepackage{booktabs}
\usepackage{float}
\usepackage{array}
\usepackage[english]{babel}
\usepackage{subfigure}
\usepackage{ragged2e}
\usepackage[bottom]{footmisc}
\usepackage{mathtools}
\usepackage[thinc]{esdiff}

\DeclareMathOperator{\E}{\mathbb{E}}


% \title{Introduction to Thesis "Statistical Inference for State-Space Models and Their Application to Complex Environments"}
\title{Summary of Doctoral Thesis \\ "Strategic Behavior and Economic Dynamics in ICT Markets: An Agent-Based Framework"}


\author{Zhongli Wang\thanks{PhD student affiliated with Bielefeld University (UNIBI) and Catholic University of Milan (UCSC). Email address: \url{zhongli.wang@uni-bielefeld.de}. Many thanks for the guidance of Prof. Herbert Dawid (UNIBI), Prof. Domenico Delli Gatti (UCSC) and Prof. Christiane Fuchs (UNIBI).}} 
%  This work has received funding from the European Union's Horizon 2020 research and innovation programme under the Marie Skłodowska-Curie grant agreement No 956107, "Economic Policy in Complex Environments (EPOC)".
%\date{February 2022} 

\begin{document}
% \include{chapter1}
\maketitle

\onehalfspacing
% \section{Introduction}

\begin{comment}
State-space models (SSMs) provide a widely used framework for analyzing dynamic systems where hidden states evolve over time based on observed variables.
\textcite{Lux2018EmpiricalModels} summarizes the framework of a general state-space model as: 
\begin{equation*}
\begin{aligned}
    x_t & = f(x_{t-1}, \epsilon_t) \\
    y_t & = g(x_t, \eta_t)     
\end{aligned}
\end{equation*}
where $x_t$ represents the stochastic evolution in time of a vector of states, $\epsilon_t$ represents a summary for all stochastic factors that enter into the dynamics of $x_t$, $y_t$ represents a vector of measurements, $eta_t$ represents a summary for all stochastic factors that make the vector of measurements a noisy signal of the states $x_t$.
While SSMs are traditionally applied in econometrics and control theory, this dissertation explores their broader applicability in economic modeling. 
\vspace{5mm}

Specifically, although the primary models developed in this work are small-scale agent-based models (ABMs) rather than classical SSMs, they share key methodological features. 
\textcite{Tesfatsion2006ChapterTheory} illustrates that the agent-based modeling specifies the initial state of an economic system by specifying each
agent’s initial data and behavioral methods and the degree of accessibility of these data
and methods to other agents, which shares similar features with general SSMs.
\textcite{Lux2018EmpiricalModels} further summarizes that the estimation of ABMs could profit substantially from the methods for state-space models or slightly more general latent variable models. Moreover, many ABMs, including the models in this thesis, to a large extent satisfy the general framework of a general SSM mentioned above. 
In particular, ABMs capture system dynamics through agent interactions, which collectively define an evolving state over time. This conceptual similarity allows for the application of statistical inference techniques inspired by SSMs, such as parameter estimation and model validation, to ABMs. By adopting this perspective, the dissertation aims to bridge elements of state-space analysis and agent-based modeling, while recognizing the fundamental differences between the two approaches.
\vspace{5mm}

Based on the "bottom-up" nature, the ABMs consider the microfoundations of the economic system and have the potential to show emergent behavior in a complex world. Abundant literature suggests potential of this type of models, particularly in macroeconomics \parencite{Axtell2022Agent-BasedFuture} and digital economics \parencite{DawidDelliGatti2018, Dawid2021IntroductionModel}.
From theory to data, researches related to calibrating ABMs include frequentist approaches, e.g. method of moments, maximum likelihood, and recent Bayesian inference methods, e.g. Bayesian Monte-Carlo and Bayesian kernels. Other advanced methods including neural networks have also been applied \parencite{Grazzini2015EstimationDistance, Lux2018EmpiricalModels, DelliGatti2020RisingModels, Platt2022BayesianNetworks}. 
Nevertheless, it is important to acknowledge that at the meso-level of economics - particularly when examining competitive dynamics within specific industries - existing approaches may overlook critical complexities. 
The Information and Communications Technology (ICT) industry provides a compelling example, characterized by intricate dynamics such as platform competition, data restrictions, user experience, and capacity constraints. These complexities have not been sufficiently addressed in prior literature, neither at the theoretical level nor within empirical studies that primarily emphasize structural parameter estimation.


\vspace{5mm}
\end{comment}

The dissertation consists of three main chapters. The first two chapters are theoretical, developing models to analyze firm strategies and market dynamics. While both focus on theoretical foundations, they differ in scope: 
Chapter 1 establishes the core theoretical framework of an IO-related model, examining how data restriction policies influence firm profits and social welfare.
Chapter 2 extends the analysis by incorporating consumer networks and consumer congestion effects, exploring their impact on market outcomes and firm competition.
The final chapter (Chapter 3) is empirical, applying the theoretical foundations of the first two chapters to real-world data. By integrating firm-level data, it assesses the practical relevance and applicability of the proposed models, bridging the gap between theory and empirical validation.
\vspace{5mm}


Chapter 1 examines the impact of data restrictions on the software market, focusing on how varying levels of regulation influence the geographical distribution of economic activities over time. Unlike traditional goods, software products are inherently intangible, rely heavily on information systems, and operate under complex licensing mechanisms. Their pricing and availability differ across regions due to variations in local data regulations.
To investigate these dynamics, this chapter adopts a two-fold approach, combining game theory and agent-based modeling to simulate interactions among key market participants - software owners and developers, consumers, and governments. The proposed models capture firms’ strategic behavior, consumers’ decision-making processes, and the dynamic pricing mechanisms shaped by evolving regulatory landscapes.
Findings from computational experiments highlight that data restrictions play a crucial role in shaping price formation and market structures. Stricter regulations can distort prices, alter competitive strategies, and reshape market dynamics, underscoring the need for a balanced regulatory framework that supports both market efficiency and consumer protection.
\vspace{5mm}

Chapter 2 assess the interplay between firm capacity selection, network effects, and market competition in the ICT industry. Traditional industrial organization theory emphasizes R\&D investment as a primary driver of firm profitability through innovation. However, in varied market structures, this relationship becomes more complex due to firms' service capacity constraints. Exceeding these limits can degrade product quality, reduce consumer retention, and ultimately impact profitability. Additionally, consumer behavior - strongly shaped by digital community interactions - plays a crucial role in determining the perceived value of R\&D-driven quality improvements. The model in this chapter captures firms' investment strategies, the decision-making processes of consumers, and the broader influence of consumer networks on market outcomes.
The findings provide new insights into the interdependencies between service capacity, R\&D investment effectiveness, and consumer engagement. By investigating how firms navigate these challenges, this chapter contributes to a deeper understanding of strategic capacity selection in competitive environments and the role of network effects in shaping firm performance. The results highlight the importance of balancing capacity constraints, investment efficiency, and consumer relationships to sustain long-term profitability in the ICT sector.


\vspace{5mm}

Chapter 3 explores advancements in parameter estimation and calibration in agent-based models (ABMs). While ABMs have significant potential to uncover emergent phenomena and inform policy decisions, the estimation and calibration of model parameters remains a major challenge. This chapter addresses this issue by assessing the model setup introduced in Chapter 2, where strategic firm behavior and network effects play a central role.
To enhance model validity, this chapter employs state-of-the-art estimation techniques, including the Simulated Method of Moments (SMM) and Bayesian Monte Carlo methods, integrating financial data from major U.S. telecommunication firms (AT\&T and T-Mobile) to calibrate parameters and validate the model. By systematically comparing these methodologies, this chapter evaluates their suitability, effectiveness, and practical implications for agent-based economic modeling.
Beyond methodological advancements, this chapter contributes to the broader discourse on improving the validity and applicability of computational economic models. The findings provide key insights into overcoming industry-specific challenges, enhancing model reliability, and bridging the gap between theoretical modeling and real-world applications.

% \section*{Acknowledgments} 

% This work has received funding from the European Union's Horizon 2020 research and innovation programme under the Marie Skłodowska-Curie grant agreement No 956107, "Economic Policy in Complex Environments (EPOC)".


\newpage
\printbibliography

\end{document}

State-space models ... Example of ABMs, philosophy

This thesis handles the above mentioned economics and statistical problems ... 
\vspace{5mm}

From the aspect of IO, ABMs 

There is abundant evidence that ABMs can show the feature

connection with IO literature, applications
\vspace{5mm}

especially digital economics, applications
\vspace{5mm}

Nevertheless, it is worthwhile to mention that on the meso level of economics, for specific examples in competition dynamics within a single industry, the complexity may be overlooked. The Information and Communications Technology (ICT) industry is a good example, since it involves many complex notions including platform competition, data restrictions, user experience, usage capacities, etc. which are not specified in previous literature, either on the theoretical level or on the empirical level (which focuses on structural parameter estimation).



There are three main chapters, two are theoretical, the last one is empirical and further suggests the application of the theoretical background of the first two chapters.

The dissertation consists of three main chapters. The first two chapters are theoretical, developing models to analyze firm strategies and market dynamics. 
The difference in first two chapters is that the Chapter 1 develops the basic theoretical structure of the IO-related model and assess the influence of the data restriction policies on the firm profits and social welfare, while Chapter 2 evaluates the consumer network and the consumer congestion and aims to analyze the possible market outcomes in terms of firm competition.  
The final chapter is empirical, applying the theoretical foundations established in the first two chapters to real-world data, demonstrating their practical relevance and applicability.


Chapter 1 studies the impact of data restrictions on the software market. It aims to answer the question of how the varied levels of data restrictions shape the geographical distribution of economic activities over time. Different from IO models handling traditional goods, software products rely heavily on information systems, are inherently intangible, and are governed by complex licensing mechanisms. Their pricing and availability vary across regions due to differences in local data regulations. 
To explore this, Chapter 1 employs a two-fold approach considering both game theory and agent-based modeling, simulating interactions among key market participants - software owners and developers, consumers, and governments. The proposed models capture the strategic behavior of software firms, the decision-making processes of diverse consumer groups, and the dynamic pricing mechanisms influenced by evolving regulatory landscapes.
Findings from the computational experiments reveal that the degree of data restrictions plays a pivotal role in shaping price formation and market structures. Stricter data regulations can lead to price distortions, altered market strategies, and shifts in competitive dynamics. 

Chapter 3 investigates advancing parameter estimation and calibration in ABMs. Despite the potential of ABMs to uncover emergent phenomena and inform policy decisions, the estimation and calibration of model parameters remains a critical challenge. This chapter addresses the issue above by assessing the model set up in Chapter 2, where strategic firm behavior and network effects play a central role.
To enhance model validity, this chapter employs state-of-the-art estimation techniques, including the Simulated Method of Moments (SMM) and Bayesian Monte Carlo methods, integrating financial data from U.S. telecommunication firms (AT\&T and T-Mobile) to calibrate parameters and validate the model. By systematically comparing these methodologies, this chapter assesses their suitability, effectiveness, and practical implications for agent-based economic modeling.
This chapter also advances the methodological rigor of ABM parameterization and contributes to the broader discourse on improving the validity and applicability of computational economic models. 
The findings offer key insights into overcoming industry-specific challenges, enhancing model reliability, and bridging the gap between theoretical modeling and real-world applications.

